\section{Results}
\label{sec:results}

All numbers are over three seeds, reported as mean $\pm$ standard deviation. A
run that fails to reach $95\%$ test accuracy within $30\,000$ steps is recorded
as a failure and excluded from averages rather than assigned the step budget; a
bracketed fraction indicates how many seeds generalised where not all did.
Ablations use seed $0$ only and are reported separately in
Section~\ref{sec:res-ablations}.

\subsection{Every architecture memorises immediately; they differ in when they stop}

Figure~\ref{fig:grokking} shows training and test accuracy for all seven layer
types at $N = 256$ composed experts. The qualitative picture is the same
everywhere: training accuracy saturates within a few hundred steps, and test
accuracy follows after a delay that ranges from a few hundred steps to never.
The delay is what the architectures disagree about.

\begin{figure}[ht]
\centering
\includegraphics[width=\textwidth]{fig1_grokking_curves.pdf}
\caption{Training (left) and test (right) accuracy at $N = 256$ experts,
seed 0. Every architecture fits the training set quickly; the gap to
generalisation varies by more than an order of magnitude, and the unfactorised
mixture (orange) never closes it.}
\label{fig:grokking}
\end{figure}

Table~\ref{tab:headline} gives the grokking step at a fixed expert count. The
ordering spans a factor of $27$ between the fastest factorisation and the
slowest, with the unfactorised mixture failing on all three seeds.

\begin{table}[ht]
\centering
\caption{Steps to $95\%$ test accuracy at $N = 256$ composed experts, mean
$\pm$ sd over three seeds. The dense MLP has no experts and is shown for
reference.}
\label{tab:headline}
\begin{tabular}{lrrr}
\toprule
Architecture & Parameters & Grokking step & Final test acc. \\
\midrule
Tucker            &  287\,968 & $\mathbf{600 \pm 100}$ & $1.000$ \\
Dense MLP         &   34\,113 & $1667 \pm 58$          & $0.996$ \\
Tensor-Train      &   78\,480 & $2033 \pm 58$          & $1.000$ \\
CP                &   61\,952 & $2267 \pm 289$         & $1.000$ \\
MONET-HD          &   59\,792 & $5333 \pm 1986$        & $0.998$ \\
MONET-VD          &   59\,792 & $16400 \pm 6861$       & $0.986$ \\
Unfactorised MoE  &  739\,712 & never (0/3)            & $0.001$ \\
\bottomrule
\end{tabular}
\end{table}

Two observations do not follow from the ordering itself. The unfactorised
mixture holds the largest parameter budget of the seven and is the only one
that fails, so capacity is not what is being measured. And the spread across
seeds differs by two orders of magnitude between architectures: Tucker's seeds
grok at $500$, $600$ and $700$ steps, while MONET-VD's grok at $8600$, $19100$
and $21500$. Reliability under reinitialisation is itself a function of the
factorisation; we return to this in Section~\ref{sec:res-stability}.

\subsection{The expert-count effect has no single sign}
\label{sec:res-scaling}

Table~\ref{tab:scaling} gives the grokking step as the expert count grows from
$16$ to $4096$ for every architecture.

\begin{table}[ht]
\centering
\caption{Steps to $95\%$ test accuracy against expert count, mean $\pm$ sd over
three seeds. ``fail'' means no seed reached the threshold; a bracketed fraction
gives the number that did where some did not. The unfactorised layer is not run
beyond $N = 1024$.}
\label{tab:scaling}
\footnotesize
\begin{tabular}{lrrrrrr}
\toprule
$N$ & Tucker & CP & Tensor-Train & MONET-HD & MONET-VD & Unfactorised \\
\midrule
   16 & $967 \pm 208$ & $7900$\,[1/3]  & $1300 \pm 361$ & $2967 \pm 351$  & fail & $2733 \pm 252$ \\
   64 & $900 \pm 436$ & $2967 \pm 404$ & $1767 \pm 907$ & $5233 \pm 702$  & fail & $7733 \pm 1762$ \\
  256 & $600 \pm 100$ & $2267 \pm 289$ & $2033 \pm 58$  & $5333 \pm 1986$ & $16400 \pm 6861$ & fail \\
 1024 & $500 \pm 100$ & $2067 \pm 58$  & $2500 \pm 436$ & $8500 \pm 656$  & fail & fail \\
 4096 & $533 \pm 58$  & $2000 \pm 100$ & $4567 \pm 709$ & $21850 \pm 6010$\,[2/3] & fail & --- \\
\midrule
trend & $\downarrow$ & $\downarrow$ & $\uparrow$ & $\uparrow$ & --- & collapse \\
\bottomrule
\end{tabular}
\end{table}

Adding experts helps two architectures, hurts two, and destroys one. Tucker
improves by a factor of $1.9$ from $N = 16$ to $N = 1024$ and still generalises
on every seed at $N = 4096$. CP improves by $4.0\times$, and the shape of that
improvement is the more telling half of it: at $N = 16$ only one seed in three
generalises at all, so for CP a small expert count is the failure regime.
Tensor-Train moves the other way, degrading by $3.5\times$, and MONET-HD by
$7.4\times$, with one seed in three no longer generalising at $N = 4096$. The
unfactorised mixture generalises at $N = 16$ and $N = 64$ and then fails on
every seed at $N = 256$, reaching training accuracy $1.000$ with test
accuracies $0.0004$, $0.0000$ and $0.0027$.

MONET-VD stands apart: it fails at $N = 16$, $64$, $1024$ and $4096$, and
generalises only at $N = 256$. This is not a trend with an exception but a
single window in an otherwise failing architecture, and we have no explanation
for it.

\begin{figure}[ht]
\centering
\includegraphics[width=0.85\textwidth]{fig2_grok_vs_experts.pdf}
\caption{Steps to generalisation against the number of composed experts.
Crosses mark configurations where no seed reached the threshold. The curves do
not share a direction: Tucker and CP slope down, Tensor-Train and MONET-HD
slope up, and the unfactorised mixture terminates.}
\label{fig:scaling}
\end{figure}

The claim that more experts is better is therefore not a property of mixtures
of experts. Nor --- and this is where our initial expectation was wrong --- is
it settled by whether the expert tensor is factorised at all. Tensor-Train shares a basis as thoroughly as Tucker does, and keeps it as clean at the two middle expert counts
(Table~\ref{tab:basis-scaling}), yet its grokking step rises with scale while
Tucker's falls. What the results establish is weaker and more specific: the
sign of the effect is a property of the particular factorisation, and a claim
of the form ``more experts is better'' has to be made one architecture at a
time.

\subsection{What survives scaling is the basis, not the router}

The task supplies an independent view of the same models: since the
generalising algorithm must represent operands through Fourier frequencies, we
can ask which frequencies each layer uses, and how that changes as experts are
added.

Figure~\ref{fig:spectra} shows the pooled power spectrum of the shared basis at
$N = 256$ for seed $0$. Averaged over three seeds, the dense MLP spreads its
representation almost evenly, using $38.1$ of the $56$ available frequencies;
MONET-HD concentrates on $7.6$; the Tucker layer solves the same task with $1.7$. A layer with $287\,968$ parameters represents the task in a
two-dimensional frequency subspace, while a layer with $34\,113$ parameters
spreads it over nearly forty.

\begin{figure}[ht]
\centering
\includegraphics[width=\textwidth]{fig5_spectra.pdf}
\caption{Pooled Fourier power of the shared basis as a function of the operand,
at $N = 256$, seed $0$. The effective frequency count is the participation ratio
of the pooled spectrum; three-seed means are given in
Table~\ref{tab:basis-scaling}.}
\label{fig:spectra}
\end{figure}

Table~\ref{tab:basis-scaling} tracks this along the expert-count axis, and it
separates the architectures differently from Table~\ref{tab:scaling}. Tucker
and Tensor-Train both hold a narrow, clean basis as the layer grows: Tucker
between $1.5$ and $1.8$ effective frequencies across a $256\times$ range, and
Tensor-Train at $1.1$ and $1.2$ at $N = 256$ and $N = 1024$, with purity above
$0.94$ at those two counts.
CP broadens steadily ($14.1 \to 35.6$) while its purity nonetheless rises, and
the product-key layers deteriorate on both measures at once --- MONET-HD from
purity $0.891$ and $8.5$ frequencies to $0.150$ and $29.3$, MONET-VD from
$0.881$ and $15.9$ to $0.276$ and $53.4$, against a white-noise purity
reference of $0.082$.

\begin{table}[ht]
\centering
\caption{Effective number of frequencies in the shared basis (top) and
norm-weighted spectral purity (bottom), mean over three seeds. The task has
$56$ frequencies in total; random weights score purity $0.082$.}
\label{tab:basis-scaling}
\footnotesize
\begin{tabular}{lrrrrrr}
\toprule
$N$ & Tucker & CP & Tensor-Train & MONET-HD & MONET-VD & Unfactorised \\
\midrule
\multicolumn{7}{l}{\emph{effective frequencies}} \\
   16 & $1.8$ & $14.1$ & $4.1$  & $8.5$  & $15.9$ & $36.4$ \\
   64 & $1.5$ & $19.2$ & $2.9$  & $8.5$  & $14.5$ & $28.4$ \\
  256 & $1.7$ & $21.7$ & $1.1$  & $7.6$  & $15.9$ & $53.6$ \\
 1024 & $1.7$ & $31.9$ & $1.2$  & $15.0$ & $39.8$ & $55.5$ \\
 4096 & $1.6$ & $35.6$ & $18.6$ & $29.3$ & $53.4$ & --- \\
\midrule
\multicolumn{7}{l}{\emph{basis purity}} \\
   16 & $0.865$ & $0.611$ & $0.713$ & $0.891$ & $0.881$ & $0.824$ \\
   64 & $0.940$ & $0.780$ & $0.732$ & $0.811$ & $0.864$ & $0.796$ \\
  256 & $0.914$ & $0.867$ & $0.948$ & $0.691$ & $0.834$ & $0.138$ \\
 1024 & $0.953$ & $0.879$ & $0.963$ & $0.242$ & $0.606$ & $0.091$ \\
 4096 & $0.951$ & $0.906$ & $0.682$ & $0.150$ & $0.276$ & --- \\
\bottomrule
\end{tabular}
\end{table}

Tensor-Train's apparent collapse at $N = 4096$ is an averaging artefact worth
naming: two of its three seeds hold purity $0.976$ and $0.978$ at $1.11$ and
$1.16$ frequencies, while the third drops to $0.092$ at $53.5$ frequencies ---
and reaches $100\%$ test accuracy anyway, by some route we have not identified.
The mean of $0.682$ describes neither.

The configurations that fail to generalise are the endpoint of this trend
rather than a separate phenomenon. The unfactorised mixture reaches purity
$0.138$ with $53.6$ effective frequencies at $N = 256$, and $0.091$ with $55.5$
of the available $56$ at $N = 1024$: a representation indistinguishable from
noise, produced by a network that fits its training set perfectly.

Across the mixture family, narrower frequency budgets accompany earlier
generalisation, but only loosely: over the $66$ generalising mixture runs the rank correlation between effective frequency count and grokking step is $\rho = 0.50$
(Figure~\ref{fig:freqgrok}). These runs share seeds and their
architectures are nested, so the coefficient summarises the scatter rather than
testing a hypothesis. Two systematic exceptions are visible in it. The dense MLP
uses the widest frequency budget of any model here and still generalises faster
than every product-key variant, and CP broadens its basis steadily with scale
while its grokking step falls. Frequency sparsity is therefore neither necessary
nor sufficient for fast generalisation, and we describe the relationship as a
regularity within the family rather than as a mechanism.

\begin{figure}[ht]
\centering
\includegraphics[width=0.8\textwidth]{fig3_freqs_vs_grok.pdf}
\caption{Effective number of frequencies in the shared basis against steps to
generalisation, one point per generalising run. Note the dense MLP (blue), which combines the widest frequency budget with one of the lowest grokking steps.}
\label{fig:freqgrok}
\end{figure}

\paragraph{The router empties out.}
\label{sec:res-router-consequence} Applying the same analysis to the router
gives the most uniform result in this study. Under one-hot inputs the router
logits are a learned lookup table indexed by the operand, so their spectrum is
as directly interpretable as the basis's. Table~\ref{tab:router-scaling} shows
what happens to it.

\begin{table}[ht]
\centering
\caption{Norm-weighted spectral purity of the router keys against expert count,
mean over three seeds. Values at or below the white-noise reference of $0.082$
are in bold. Every architecture ends at or near the floor, though the descent is not monotone for the product-key layers.}
\label{tab:router-scaling}
\footnotesize
\begin{tabular}{lrrrrrr}
\toprule
$N$ & Tucker & CP & Tensor-Train & MONET-HD & MONET-VD & Unfactorised \\
\midrule
   16 & $0.622$ & $0.575$ & $0.607$ & $0.204$ & $0.336$ & $0.552$ \\
   64 & $0.525$ & $0.258$ & $0.458$ & $0.215$ & $0.632$ & $0.570$ \\
  256 & $0.143$ & $0.143$ & $0.177$ & $0.163$ & $0.383$ & $\mathbf{0.089}$ \\
 1024 & $\mathbf{0.072}$ & $\mathbf{0.077}$ & $\mathbf{0.082}$ & $0.192$ & $0.194$ & $\mathbf{0.074}$ \\
 4096 & $\mathbf{0.048}$ & $\mathbf{0.056}$ & $\mathbf{0.059}$ & $0.093$ & $\mathbf{0.076}$ & --- \\
\bottomrule
\end{tabular}
\end{table}

At $N = 16$ every router carries visible structure, scoring between $0.20$ and
$0.62$ against a noise floor of $0.082$. By the largest expert count at which each architecture is run, every router has
fallen to between $0.048$ and $0.093$: at or below the noise floor for five of
the six, and within $14\%$ of it for MONET-HD. The decay holds across all six architectures --- the ones that generalise faster with scale and the ones
that generalise slower, the ones whose bases stay clean and the ones whose
bases decay alongside. Whatever else the choice of factorisation decides, it
does not decide this: adding experts empties the router of interpretable
structure.

The contrast with the basis is sharpest where the basis survives. A Tucker
layer at $N = 4096$ has basis purity $0.951$ and router purity $0.048$: the
computation sits entirely in the shared factors, and the the routing weights carry no more frequency structure than random numbers. Figure~\ref{fig:purity} compares
the two components across architectures, and Figure~\ref{fig:tables} shows them
side by side for a single layer. A complementary measure --- the overlap
between the pooled router spectrum and the pooled basis spectrum, computed on
seed $0$ --- gives the same ordering, with Tucker layers at $0.01$--$0.10$
against $0.46$--$0.78$ for the product-key layers and $0.93$ for the
unfactorised mixture that fails.

\paragraph{The decay is not what empties the router.} Weight decay is applied
to the router keys throughout, which raises the possibility that the collapse
in Table~\ref{tab:router-scaling} is regularisation rather than scale. The
runs that exempt the keys settle it. Exempting them leaves the keys three to
five times larger in root-mean-square magnitude --- $0.557$ against $0.152$
for MONET-HD at $N = 4096$, $0.905$ against $0.336$ at $N = 256$, and $0.270$
against $0.077$ for Tucker at $N = 4096$ --- and does not restore their
spectrum: purity is $0.069$ against $0.093$, $0.119$ against $0.163$, and
$0.045$ against $0.048$ respectively. In every pair the unregularised router
is larger and no better structured. Magnitude and frequency content are
independent here, and it is the second that scaling removes. Note also that
the reference value $0.082$ is the mean purity of random tables of the same
shape rather than a lower bound, so measured values slightly below it are
unremarkable.

\begin{figure}[ht]
\centering
\includegraphics[width=\textwidth]{fig4_purity_basis_vs_router.pdf}
\caption{Norm-weighted spectral purity of the shared expert basis (left)
and of the router keys (right) against expert count, averaged over three
seeds. The dotted line is the purity of a white-noise table ($0.082$).
The dense MLP has neither a basis nor a router and is therefore absent.
The Tucker and CP bases stay clean across the whole sweep, and Tensor-Train's
does except for one seed at $N = 4096$; every router decays to the noise floor.}
\label{fig:purity}
\end{figure}

\begin{figure}[ht]
\centering
\includegraphics[width=\textwidth]{fig6_tables.pdf}
\caption{Shared basis (left) and router keys (right) of a Tucker layer as
functions of the first operand, rows sorted by dominant frequency and
normalised individually, with pooled spectra below. The basis is periodic and
concentrated; the router is not.}
\label{fig:tables}
\end{figure}

This is the observation with the most direct consequence for practice, and the
next section shows how far it reaches. MONET identifies specialised experts by
routing-score skew, and operates at $N = 262\,144$ --- far beyond the point
where, on this task, the router's frequency content has already fallen to
noise. What such a criterion can still recover at that scale, and what it
cannot, is the subject of Section~\ref{sec:res-mixed}.

\subsection{A positive control: what routing scores do recover}
\label{sec:res-mixed}

Every measurement so far is frequency-based, because on single-task modular
addition there is no partition of the input to be specialised towards. The
mixed dataset of Section~\ref{sec:method} supplies one: addition and
multiplication examples over the same operand grid, distinguished by a task
token appended to the input. Here class-based criteria apply, so MONET's own
procedure can be run end to end and checked against ablation --- on a task
where the partition is known exactly. It serves as a positive control: if the
procedure finds nothing here, the negative results above would be evidence
about our measurements rather than about the models.

Both architectures learn the mixed task, and the expert-count trends of
Section~\ref{sec:res-scaling} survive it in weakened form. Tucker generalises
in $1700$, $1800$ and $2100$ steps at $N = 256$, $1024$ and $4096$, so its
improvement with expert count disappears; MONET-HD fails at $N = 256$ and
takes $8800$ and $17300$ steps at the two larger counts, degrading about twice
as steeply as on the single task. Adding a second subtask slows Tucker at every expert count and narrows MONET-HD's window from below.

\paragraph{Procedure.} We follow MONET: a unit is assigned to a subtask when
its mean routing score on that subtask exceeds its score on the other by a
factor of two, and the assignment is then checked by ablation. Two adaptations
are forced by product-key composition. First, a composed expert $(i,j)$ cannot
be ablated individually --- the layer never materialises it, contracting $g^{(1)}$
and $g^{(2)}$ separately --- so we score and ablate the $2Hn$ \emph{factors}
instead, which are the objects the parameters actually correspond to. Second,
ablation sets the factor's routing weight to zero without renormalising the
remainder, so a removed factor contributes nothing, including through its bias
terms. We report accuracy separately on each subtask, and against a control
that ablates the same number of factors chosen uniformly at random, averaged
over twenty draws.

\begin{table}[ht]
\centering
\caption{Applying MONET's specialisation criterion on the mixed task, seed $0$.
``Flagged'' counts the factors assigned to addition, out of $2Hn$. $\Delta$ is
the drop in accuracy on each subtask when those factors are ablated, and
selectivity is $\Delta_{\text{add}} - \Delta_{\text{mul}}$. The control ablates
the same number of factors at random, averaged over twenty draws. MONET-HD at
$N = 256$ did not generalise and is excluded.}
\label{tab:mixed}
\footnotesize
\begin{tabular}{lrrrrrr}
\toprule
 & & \multicolumn{3}{c}{criterion} & \multicolumn{2}{c}{random control} \\
\cmidrule(lr){3-5}\cmidrule(lr){6-7}
Architecture & $N$ & flagged & $\Delta_{\text{add}}$ & selectivity
& $\Delta_{\text{add}}$ & selectivity \\
\midrule
MONET-HD & $1024$ & $33/128$ & $0.981$ & $\mathbf{0.983}$ & $0.601$ & $0.121$ \\
MONET-HD & $4096$ & $51/256$ & $0.965$ & $\mathbf{0.969}$ & $0.269$ & $-0.237$ \\
Tucker   & $256$  & $31/64$  & $0.987$ & $\mathbf{0.987}$ & $0.233$ & $0.211$ \\
Tucker   & $1024$ & $51/128$ & $0.993$ & $\mathbf{0.993}$ & $0.021$ & $-0.058$ \\
Tucker   & $4096$ & $67/256$ & $0.991$ & $\mathbf{0.991}$ & $0.005$ & $0.003$ \\
\bottomrule
\end{tabular}
\end{table}

\paragraph{The criterion passes, and passes at every scale.} Table~\ref{tab:mixed}
is unambiguous. Ablating the factors that routing scores assign to addition
removes addition entirely --- accuracy falls by $0.97$ to $0.99$ --- and leaves
multiplication untouched, with $\Delta_{\text{mul}}$ between $-0.004$ and
$0.000$ in every configuration. The random control removes the same number of
factors and damages both subtasks indiscriminately, with selectivity ranging
from $0.21$ down to $-0.24$. The effect does not weaken with expert count: at
$N = 4096$, where the same routers score $0.048$ and $0.093$ on spectral
purity against a white-noise floor of $0.082$, selectivity is still $0.99$.

One feature of the table is a check on the procedure rather than a result.
No factor on the first side of the router is ever flagged. With $d_{\text{in}} = 228$ the router's split puts coordinates $0$--$113$ on the
first side --- the whole of $\mathrm{onehot}(a)$ together with the single
coordinate for $b = 0$ --- and everything else, including both task-token
coordinates, on the second, so the first side cannot
see which subtask an example belongs to and correctly reports no
specialisation. The method is not producing noise; it finds structure exactly
where structure exists.

\paragraph{What this does to the router result.} The two halves of the study
now say something sharper than either does alone. The same routers that have
lost their frequency content by $N = 4096$ still carry the subtask partition
perfectly, and routing-score attribution recovers it with near-perfect
selectivity. A router is therefore not emptied by scale in any absolute sense.
What it loses is the structure that implements the algorithm; what it keeps is
the structure the input already announces.

This is a worse position for routing-score interpretability than a simple
failure would be. The criterion and the ablation agree here because both read
the same partition, so their agreement validates the partition rather than the
mechanism, and a method that checks one against the other cannot tell the
difference. MONET's domains --- Python, chemistry, toxic text --- are partitions
of the first kind, largely readable from the tokens themselves. A routing-based
criterion will find experts that track them, the ablation check will confirm
it, and none of that is evidence about how the model computes anything.

\subsection{Ablations}
\label{sec:res-ablations}

Four ablations at seed $0$ --- router-key decay at both $N = 256$ and
$N = 4096$, the other three at $N = 256$ --- isolate mechanisms that
would otherwise remain conjectural. Results are in Table~\ref{tab:ablations};
the baseline column repeats the corresponding main-grid run for the same seed.

\begin{table}[ht]
\centering
\caption{Ablations. ``fail'' gives the final test accuracy in parentheses. The
baseline configuration is a natural input split, $\lambda = 10^{-3}$, router
keys decayed, $80\%$ training fraction. Rows marked $\dagger$ are over three
seeds, reported as mean $\pm$ sd; all other entries are seed $0$.}
\label{tab:ablations}
\footnotesize
\begin{tabular}{llrrr}
\toprule
Ablation & Architecture & \multicolumn{3}{c}{Grokking step} \\
\midrule
\emph{Router-key decay} & & decayed & exempt & \\
& Tucker, $N{=}256$       & $600$  & $600$ & \\
& Tucker, $N{=}4096$      & $600$  & $600$ & \\
& MONET-HD, $N{=}256^\dagger$  & $5333 \pm 1986$ & $5800 \pm 929$ & \\
& MONET-HD, $N{=}4096^\dagger$ & $21850 \pm 6010$\,[2/3] & fail $(0.006$--$0.090)$ & \\
& Unfactorised, $N{=}256$ & fail $(0.000)$ & fail $(0.007)$ & \\
\midrule
\emph{Router input split} & & natural & interleaved & random \\
& Tucker, $N{=}256$   & $600$  & $5100$ & $1700$ \\
& MONET-HD, $N{=}256$ & $7600$ & fail $(0.823)$ & fail $(0.822)$ \\
\midrule
\emph{Auxiliary loss $\lambda$} & & $0$ & $10^{-3}$ & $10^{-2}$ \\
& Tucker, $N{=}256$   & $500$  & $600$  & $800$ \\
& MONET-HD, $N{=}256$ & $6500$ & $7600$ & $3000$ \\
\midrule
\emph{Training fraction} & & $0.4$ & $0.6$ & $0.8$ \\
& Tucker, $N{=}256$   & $6300$ & $800$  & $600$ \\
& MONET-HD, $N{=}256$ & fail $(0.004)$ & $8200$ & $7600$ \\
\bottomrule
\end{tabular}
\end{table}

\paragraph{Router-key decay matters where memorisation is cheap.} Under one-hot
inputs the router keys form a table of $8nP$ parameters addressed directly by the operand value --- $14\,464$ of them at $N = 256$ and $57\,856$ at $N = 4096$.
Exempting them from weight decay, as is conventional for router parameters,
leaves that table unregularised. At the $80\%$ training fraction used throughout, the cost appears only where
the table is largest, and it is decisive there. Over three seeds at
$N = 4096$, exempting the keys costs generalisation on every one, with final
test accuracies of $0.006$, $0.013$ and $0.090$, against two successes in
three for the decayed setting; at $N = 256$ the two settings are
indistinguishable, $5333 \pm 1986$ against $5800 \pm 929$. For Tucker the ablation makes no difference at either expert count.

The effect is far larger when data is scarcer. During protocol selection
(Section~\ref{sec:protocol}) we compared the two settings at a $50\%$ training
fraction with $N = 256$: with the keys exempt, test accuracy was $0.006$; with
the keys decayed, the layer grokked at step $15\,000$ and reached $0.977$.
Both the data fraction and the expert count control how attractive the
memorisation shortcut is, and the convention of exempting router parameters is
safe only while it stays unattractive --- which is not something one can assume
without checking.

\paragraph{The natural input split matters, and unequally.} When the router's
two halves are the two operands, the factorisation of the router matches the
factorisation of the task. Permuting the input destroys that alignment.
MONET-HD then fails under both permutations, reaching only $0.82$ test
accuracy; Tucker survives both, slowed by $8.5\times$ under interleaving and
$2.8\times$ under a random permutation. Part of the advantage we report is
therefore due to structural alignment rather than to factorisation as such, and
the effect is large enough to state plainly rather than file under caveats. It
is also informative that the two architectures depend on it to different
degrees: alignment is what lets the product-key layer work here at all, while
for Tucker it is an accelerant. We have no account of why interleaving costs Tucker more than a random
permutation does.

\paragraph{The auxiliary losses do little.} MONET's uniformity and ambiguity
terms move Tucker's grokking step from $500$ at $\lambda = 0$ to $600$ and
$800$ as $\lambda$ rises --- a mild cost. For MONET-HD the three values give $6500$, $7600$ and $3000$: not monotone, and
the drop at $\lambda = 10^{-2}$ is $2.3$ standard deviations of that
configuration's seed spread, which a single seed cannot separate from noise. We
record it as unresolved rather than as an effect. Since these terms were calibrated for language
modelling, where cross-entropy never approaches zero, their near-neutrality on
an algorithmic task is unsurprising. What is worth recording is that they are
not a prerequisite for anything we observe: the frequency structure and the
router decay both appear at $\lambda = 0$.

\paragraph{Data efficiency follows the same ordering.} At a $40\%$ training
fraction Tucker still generalises, in $6300$ steps, while MONET-HD does not
generalise at all ($0.004$ test accuracy). At $60\%$ Tucker takes $800$ steps
and MONET-HD $8200$. The gap widens as data is withdrawn, which is what one
expects if the shared basis acts as a prior rather than as extra capacity.

\subsection{The rank effect has no single sign either}
\label{sec:res-rank}

The architectures that improve with expert count share a property: their
capacity along the expert modes grows with $n$. Our Tucker layer uses an
expert-mode rank of $\min(32, n)$ and CP projects onto $n$ columns, so both
widen as experts are added. Tensor-Train does not: its two expert-mode cores
are joined through a bond whose rank we held fixed at $16$ across the whole
sweep. A natural account of why Tensor-Train degrades while Tucker improves is
therefore that its added experts cannot refine the computation and can only be
used to fit the training set. Two experiments test that account, and they
disagree with each other.

\begin{table}[ht]
\centering
\caption{Rank sweep, seed $0$. Grokking step, with parameter count in
parentheses for the Tensor-Train rows. ``fail'' gives the final test accuracy.
The Tucker row freezes the expert-mode rank at $4$ so that it cannot grow with
$n$; the main-grid Tucker layer, whose rank is $\min(32, n)$, generalises in
$600$ steps at both $N = 256$ and $N = 4096$ on this seed.}
\label{tab:rank}
\footnotesize
\begin{tabular}{lccc}
\toprule
 & $N = 256$ & $N = 1024$ & $N = 4096$ \\
\midrule
\multicolumn{4}{l}{\emph{Tensor-Train}} \\
rank $8$  & $1400$ ($31$K)           & $2800$ ($47$K)            & $4300$ ($78$K) \\
rank $16$ & $2100$ ($78$K)           & $2000$ ($97$K)            & $4700$ ($135$K) \\
rank $32$ & fail $(0.866)$ ($266$K)  & $5600$ ($298$K)           & $7100$ ($360$K) \\
rank $64$ & fail $(0.022)$ ($1014$K) & fail $(0.187)$ ($1095$K)  & $16\,300$ ($1257$K) \\
\midrule
\multicolumn{4}{l}{\emph{Tucker}} \\
rank $4$, frozen & fail $(0.301)$ & $1800$ & $2400$ \\
\bottomrule
\end{tabular}
\end{table}

\paragraph{Giving Tensor-Train more rank makes it worse, not better.} The
prediction fails on both counts. The grokking step still rises with expert
count at every rank tested --- from $1400$ to $4300$ at rank $8$, from $2100$
to $4700$ at rank $16$, from $5600$ to $7100$ at rank $32$ --- so the upward
trend is not an artefact of the bottleneck we happened to fix. And raising the
rank does not help at any expert count. It helps at no expert count but one --- rank $16$ beats rank $8$ at $N = 1024$,
$2000$ steps against $2800$ --- and beyond rank $16$ it opens a failure region that widens as the rank grows: at rank $64$ and
$N = 256$ a layer holding over a million parameters fits the training set
perfectly and reaches test accuracy $0.022$.

CP separates two things the rest of this section treats together. Its units
become purer as experts are added --- purity rises from $0.611$ to $0.906$ ---
while the layer as a whole spreads over more frequencies, from $14.1$ to
$35.6$. By the per-unit criterion CP is the most monosemantic layer here; by
the layer-level criterion it is among the least sparse. Monosemanticity of
units and sparsity of the layer's frequency budget are therefore not the same
property, and the claim that a shared basis acts as a sparsity prior holds for
Tucker and Tensor-Train but not for CP.

What the sweep does track is parameter count. Rank enters the Tensor-Train
parameter count multiplicatively, so rank $64$ costs between $16$ and $33$
times rank $8$ at the same expert count, and the configurations that fail are
those where that cost is largest relative to the number of experts it buys.
This is the signature of the unfactorised mixture, which generalises at
$N = 64$ with $189\,000$ parameters and fails at $N = 256$ with $740\,000$. We
record it as a description rather than a mechanism: it orders the Tensor-Train
family and does not transfer across architectures, since a Tucker layer with
$288\,000$ parameters generalises at $N = 256$ in $600$ steps.

\paragraph{Freezing Tucker's rank removes its improvement.} The converse test
holds the Tucker expert-mode rank at $4$ so that it cannot grow with $n$. The
improvement disappears. At $N = 256$ the layer fails --- not by memorising but
by failing to fit at all, reaching training accuracy $0.432$ --- and beyond
that it degrades from $1800$ to $2400$ steps as experts are added, where the
unfrozen layer stays near $600$ throughout. On the Tucker side, then, growing
rank is exactly what converts added experts into faster generalisation.

Taken together the two experiments say something narrower than the account they
were built to test, and something that echoes
Section~\ref{sec:res-scaling}: rank growth is what lets Tucker benefit from
added experts, and the absence of rank growth is not what stops Tensor-Train.
The sign of the rank effect, like the sign of the expert-count effect, is a
property of the particular factorisation.

\paragraph{A reproducible anomaly.} Section~\ref{sec:res-stability} reports one
main-grid Tensor-Train seed that reaches perfect test accuracy at $N = 4096$
with a basis indistinguishable from noise. The rank sweep reproduces it on
demand: at rank $32$ and $N = 4096$ the layer generalises in $7100$ steps with
basis purity $0.081$ and $55.9$ effective frequencies, and at rank $64$ and
$N = 4096$ it reaches $0.998$ test accuracy with purity $0.082$. The effect is
systematic at high rank rather than a rare initialisation. We read it as a
limitation of the measurement rather than of the model: what we extract as the
shared basis of a Tensor-Train layer is its input core, and at high rank the
operand representation evidently moves elsewhere along the chain. The
Tensor-Train entries of Table~\ref{tab:basis-scaling} should be read with that
in mind.

One positive result comes out of the same runs. Tucker with rank $4$ reaches
basis purity $0.995$ and $0.993$ at $1.01$ effective frequencies for
$N = 1024$ and $N = 4096$ --- the cleanest representations anywhere in this
study, produced by layers of $31\,000$ and $60\,000$ parameters. A rank-four
expert basis solves modular addition with what is essentially a single
frequency.

\subsection{How much of the Tucker core is used}
\label{sec:res-comp}

Two facts established above point at the same object. Tucker generalises
fastest of the layers we compare, and it is also the multilinear layer whose
parameter count still rises steeply with expert count
(Table~\ref{tab:params}). Both properties live in the core, which is the only
place where the two expert modes interact. We therefore constrain that
interaction and ask how much of it the task actually uses.

Write the expert-mode slice of the core at rank $r$,
\begin{equation}
\mathcal{Z}[p, q, a, c]
  \;=\; \sum_{s=1}^{r} W[p, q, s]\, C^{(a)}[s, a]\, C^{(b)}[s, c],
\end{equation}
so that $r = 1$ makes the layer's dependence on the two routing vectors
separable, and the dense core is recovered at $r = R_a^2$, which is $256$ at
$N = 256$ and $1024$ above it. The initialisation is scaled so that the
implied core entries have the same variance at every $r$, and the rearranged
forward pass is verified against the naive sum over all $n^2$ composed experts
in the same way as the other layers.

One thing this family does \emph{not} do should be said before the numbers.
At $r = 1$ it does not reduce to product-key composition. MONET's layers apply
a nonlinearity between the two factors and are therefore not multilinear,
whereas a rank-one multilinear coupling is a strictly weaker object: a
trilinear form of rank one. The sweep varies how much the two expert modes may
interact inside a multilinear layer; it does not interpolate towards MONET.

\begin{table}[ht]
\centering
\caption{Composition-rank sweep, seed $0$. Grokking step, with parameter count
in parentheses. ``fail'' gives the final test accuracy; at ranks $1$ and $2$
the layer does not fit the training set either, reaching training accuracy
$0.23$ and $0.53$ at $N = 256$. The dense row repeats the three-seed means of
Table~\ref{tab:scaling} and the parameter counts of Table~\ref{tab:params}. The row marked $\dagger$ is over three seeds; all other rows are seed $0$.}
\label{tab:comp}
\footnotesize
\begin{tabular}{lccc}
\toprule
coupling rank & $N = 256$ & $N = 1024$ & $N = 4096$ \\
\midrule
$1$     & fail $(0.005)$ & fail $(0.007)$ & fail $(0.035)$ \\
$2$     & fail $(0.016)$ & fail $(0.029)$ & fail $(0.025)$ \\
$4$     & fail $(0.671)$ & fail $(0.475)$ & $1200$ \quad ($77$K) \\
$8$     & $800$ \quad ($34$K) & $900$ \quad ($51$K) & $800$ \quad ($82$K) \\
$16$    & $700$ \quad ($43$K) & $500$ \quad ($59$K) & $500$ \quad ($90$K) \\
$32$    & $600$ \quad ($60$K) & $500$ \quad ($77$K) & $500$ \quad ($108$K) \\
\midrule
dense$^\dagger$ & $600$ \quad ($288$K) & $500$ \quad ($1090$K) & $533$ \quad ($1121$K) \\
\bottomrule
\end{tabular}
\end{table}

\paragraph{A threshold, then nothing.} Table~\ref{tab:comp} shows a sharp
transition rather than a gradual trade-off. At ranks $1$ and $2$ the layer cannot fit the training set, let alone
generalise. That much is closer to a structural fact than to a measurement: at
$r = 1$ the layer computes a sum over heads of a linear map of the input times
a scalar function of $a$ times a scalar function of $b$, and and modular addition does not lie in that family: with $r = 1$ the
head-dependent factors enter only as scalars, so the layer computes
$s(a, b)\,\bigl(v_a(a) + v_b(b)\bigr)$ for a scalar $s$ and vectors
$v_a, v_b \in \mathbb{R}^{P}$, and the class that maximises such an expression
cannot depend on $a$ and $b$ jointly in the way $(a + b) \bmod P$ requires. What the table adds is where the threshold sits
and what happens above it. At rank $4$ it fits but
generalises only at the largest expert count. From rank $8$ upward it
generalises everywhere, and from rank $16$ it is indistinguishable from the
dense core: $700$, $500$ and $500$ steps against $600$, $500$ and $533$, all
within the seed spread that Section~\ref{sec:res-stability} reports for this
architecture.

\paragraph{The dense core is over-provisioned by an order of magnitude.} The
rank-$32$ layer at $N = 4096$ holds $107\,616$ parameters against $1\,121\,376$
for the dense core --- a factor of $10.4$ --- and generalises in $500$ steps
against $533 \pm 58$. Measured across the sweep, the parameter growth from
$N = 256$ to $N = 4096$ falls from $3.89\times$ to $1.81\times$. A coupling of rank $32$ out of a possible $1024$ is all the task uses, and rank
$8$ suffices at a cost of between a third and four fifths more steps.

This bears directly on the objection MONET raises against multilinear expert
layers (Section~\ref{sec:related}), that their parameter counts grow too
quickly with expert count to compete with product-key composition. On this
task the growth is real but almost entirely unused: it sits in core entries
that carry no signal, and constraining the coupling removes most of it at no
measurable cost in generalisation speed. The comparison between the two
families is therefore not settled by the parameter counts of the layers as
published.

\paragraph{Experts partly substitute for coupling.} At rank $4$ the layer
fails at $N = 256$ and $N = 1024$ and generalises at $N = 4096$. Where the
coupling is too poor to express the task at one expert count, adding experts
can recover it --- another instance of the window of Section~\ref{sec:res-scaling},
now with one of its two coordinates under direct control.

\subsection{Stability under reinitialisation}
\label{sec:res-stability}

A final observation emerges only with multiple seeds. At $N = 256$ the standard
deviation of the grokking step is $100$ steps for Tucker, $58$ for
Tensor-Train and the dense MLP, $289$ for CP, but $1986$ for MONET-HD and
$6861$ for MONET-VD, whose individual seeds grok at $8600$, $19100$ and
$21500$ --- a spread wider than the entire range separating the other six
architectures.

Instability also appears at the largest expert counts in architectures that are
otherwise stable, and not only in the grokking step: the Tensor-Train seed
described above, which collapses to a noise-like basis at $N = 4096$ while
still reaching perfect accuracy, has no counterpart among its siblings.
Averages at $N = 4096$ should be read with this in mind.

Product-key composition is thus not only slower but far less predictable, and a
single-seed comparison of these architectures would be seriously misleading in
either direction. This is consistent with the weak universality reported for
group-composition tasks by \citet{chughtai2023universality}, and it is why the
claims in this section are stated over seeds rather than over runs.